\section{The Fire of an Hour and the Compassion of the Saints}

Around the year 155, in the stadium of Smyrna, an eighty-six-year-old bishop named Polycarp stood before the Roman proconsul. Ordered to swear by the fortune of Caesar and revile Christ, Polycarp refused. The proconsul threatened him with wild beasts; Polycarp stood firm. Then the magistrate played his final card: ``I will cause thee to be consumed by fire, seeing thou despisest the wild beasts, if thou wilt not repent.'' Polycarp's answer has echoed through the Church for nineteen centuries:

\begin{quote}
Thou threatenest me with fire which burneth for an hour, and after a little is extinguished, but art ignorant of the fire of the coming judgment and of eternal punishment, reserved for the ungodly. But why tarriest thou? Bring forth what thou wilt.\endnote{\textit{The Martyrdom of Polycarp} 11, ANF vol.\ 1.}
\end{quote}

The chronicle notes what happened to Polycarp's face as he spoke those words: ``he was filled with confidence and joy, and his countenance was full of grace.'' The memory of eternal fire did not make the old bishop gloomy, timid, or harsh. It broke the spell of earthly terror and filled him with tranquil joy in the face of the stake.

Across the entire history of Christian holiness, the contemplation of hell has borne the same surprising fruits: steadfast courage under persecution, tender pity toward enemies, vigilance over the truth, and a burning zeal for the salvation of souls.

A generation before Polycarp's martyrdom, St.\ Ignatius of Antioch---himself on his way in chains to be torn by beasts in Rome---wrote to the Ephesians to guard them against teachers who corrupted the Gospel. If men who corrupt earthly households are punished with death, he asks, ``how much more shall this be the case with any one who corrupts by wicked doctrine the faith of God, for which Jesus Christ was crucified! Such an one becoming defiled, shall go away into everlasting fire, and so shall every one that hearkens unto him.''\endnote{St.\ Ignatius of Antioch, \textit{Epistle to the Ephesians} 16 (shorter recension), ANF vol.\ 1.} Doctrinal fidelity was for Ignatius no academic pastime; the blood of Christ had purchased the Church, and leading souls away from the truth carried eternal stakes.

A century later, during the persecutions under the Roman Empire, St.\ Cyprian of Carthage addressed a public treatise, \textit{To Demetrianus}, to a pagan magistrate who blamed Christians for every plague and drought and delighted in their torture. Cyprian does not soften the truth of the coming judgment: ``An ever-burning Gehenna will burn up the condemned, and a punishment devouring with living flames; nor will there be any source whence at any time they may have either respite or end to their torments. \dots\ The pain of punishment will then be without the fruit of penitence; weeping will be useless, and prayer ineffectual. Too late they will believe in eternal punishment who would not believe in eternal life.''\endnote{St.\ Cyprian of Carthage, \textit{To Demetrianus} 24, ANF vol.\ 5.} Yet read what Cyprian writes in the very next paragraph to the man persecuting his flock:

\begin{quote}
We repay kindness for your hatred; and for the torments and penalties which are inflicted on us, we point out to you the ways of salvation. Believe and live, and do ye who persecute us in time rejoice with us for eternity. When you have once departed thither, there is no longer any place for repentance, and no possibility of making satisfaction. Here life is either lost or saved\dots\ To him who still remains in this world no repentance is too late. The approach to God's mercy is open, and the access is easy to those who seek and apprehend the truth. Do you entreat for your sins, although it be in the very end of life, and at the setting of the sun of time; and implore God, who is the one and true God, in confession and faith of acknowledgment of Him, and pardon is granted to the man who confesses, and saving mercy is given from the divine goodness to the believer, and a passage is opened to immortality even in death itself.\endnote{St.\ Cyprian of Carthage, \textit{To Demetrianus} 25, ANF vol.\ 5.}
\end{quote}

Because hell is irrevocable after death, the mercy of Christ while life lasts becomes all the more urgent. Cyprian warns his persecutor of Gehenna for one reason only: that Demetrianus might turn before the setting of the sun and rejoice with the martyrs forever.

St.\ John Chrysostom reveals the exact same pastor's heart when he feels his congregation grow tense under his preaching on eternal fire in his ninth Homily on First Corinthians: ``Well: I know that a chill comes over you on hearing these things; but what am I to do? \dots\ The argument is irksome and pains the hearer: were it only by my own feelings, I know this. For indeed my heart is troubled and throbs; and the more I see the account of hell confirmed, the more do I tremble and shrink through fear. But it is necessary to say these things lest we fall into hell.''\endnote{St.\ John Chrysostom, \textit{Homilies on First Corinthians} 9.1--2, NPNF1 vol.\ 12.} A shepherd who hides the precipice to spare his sheep a momentary chill is not merciful; he is negligent.

When the Church moves from the patristic pulpit to the great masters of spiritual direction, the meditation on hell becomes a disciplined instrument of conversion and gratitude. In the First Week of his \textit{Spiritual Exercises}, St.\ Ignatius of Loyola places the Fifth Exercise---``a meditation on hell''---an hour before supper, applying all five interior senses to ``the length, breadth and depth of Hell.''\endnote{St.\ Ignatius of Loyola, \textit{The Spiritual Exercises}, First Week, Fifth Exercise, trans.\ Elder Mullan (New York: P.\ J.\ Kenedy \& Sons, 1914), pp.\ 44--46.} The exercitant sees in imagination the great fires and the souls as in bodies of fire, hears the wailings and blasphemies against Christ and His saints, smells smoke, sulphur, and putrid things, tastes bitter tears, sadness, and the worm of conscience, and touches how the fires burn the souls. Yet look at the grace Ignatius teaches the exercitant to beg before those five points begin:

\begin{quote}
\dots\ to ask for interior sense of the pain which the damned suffer, in order that, if, through my faults, I should forget the love of the Eternal Lord, at least the fear of the pains may help me not to come into sin.\endnote{St.\ Ignatius of Loyola, \textit{The Spiritual Exercises}, Fifth Exercise, Second Prelude, trans.\ Mullan, p.\ 45.}
\end{quote}

And look at how the Fifth Exercise closes: not in panic or gloom, but in a Colloquy directly to Christ our Lord, remembering the souls lost before His coming, during His earthly life, and after it, and giving Him heartfelt thanks ``that He has not let me fall into any of these divisions, ending my life,'' considering ``how up to now He has always had so great pity and mercy on me.''

St.\ Francis de Sales follows the same pattern in the seventh meditation of his \textit{Introduction to the Devout Life} (Part I, Chapter 15), written not for cloistered ascetics but for laypeople living in the world. After picturing hell as a dark city burning with brimstone and inhabited by citizens who cannot escape, and after weighing the pain of sense, the far greater pain of losing God's glory, and the endless night of eternity---``if a little insect in your ear, or the heat of a fever, makes one short night seem so long and tedious, how terrible will the night of eternity be''---Francis turns the heart immediately toward resolute action: ``Confess that you have often deserved it. But henceforward I will take a new course; for why should I go down into this bottomless pit? I will use such and such endeavors to avoid sin, which alone can bring me to this eternal death. Give thanks, offer, pray.''\endnote{St.\ Francis de Sales, \textit{Introduction to a Devout Life} I.15, Pustet ed., pp.\ 38--39.}

Among the saints, two approved testimonies from private revelation have entered deeply into Catholic devotion: that of St.\ Teresa of Avila, a Doctor of the Church, and that of Fatima. As private revelations, neither adds to the deposit of public Revelation closed with the Apostles, nor binds the assent of divine faith; the Church receives them as authentic graces given to stir repentance and charity.

In Chapter 32 of her \textit{Life}, written under obedience to her confessors, St.\ Teresa records that one day in prayer she found herself for a moment plunged in spirit into the place in hell which the devils kept ready for her had she continued in her former tepidity.\endnote{St.\ Teresa of Avila, \textit{The Life of St.\ Teresa of Jesus} ch.\ 32, \S\S 1--11, trans.\ Lewis, rev.\ Zimmerman, pp.\ 286--290.} She describes the low, dark, narrow passage like a furnace, the foul mire, the hollow hole in the wall where she could neither sit nor lie down nor breathe, and above all an inward anguish far worse than any bodily torture she had ever endured: ``here it is the soul itself that is tearing itself in pieces. I cannot describe that inward fire or that despair, surpassing all torments and all pain.'' Six years later, as she wrote the chapter, the natural warmth of her body still chilled at the memory.

Yet Teresa calls that moment ``one of the grandest mercies of our Lord.'' First, it destroyed her fear of earthly hardship and contradiction, making every trial of this world seem like nothing beside a single instant of that agony. Second, it poured into her soul an overwhelming gratitude to Christ her Deliverer. Third, it ignited the missionary fire that drove her to reform Carmel and found the monastery of St.\ Joseph in Avila:

\begin{quote}
It was that vision that filled me with the very great distress which I feel at the sight of so many lost souls\dots\ and also gave me the most vehement desires for the salvation of souls; for certainly I believe that, to save even one from those overwhelming torments, I would most willingly endure many deaths. \dots\ Here we know that pain ends with life at last, and that there are limits to it; yet the sight of it moves our compassion so greatly. That other pain has no ending; and I know not how we can be calm, when we see Satan carry so many souls daily away.\endnote{St.\ Teresa of Avila, \textit{The Life of St.\ Teresa of Jesus} ch.\ 32, \S 9, trans.\ Lewis, rev.\ Zimmerman, p.\ 289.}
\end{quote}

Three and a half centuries later, on 13 July 1917 at Fatima---recognized as worthy of belief by the Bishop of Leiria in 1930 and attested in Sister Lucia dos Santos's Third Memoir of 1941---the Blessed Virgin Mary briefly opened to three shepherd children a vision of hell, linking that terrible sight immediately to prayer, penance, and intercession for sinners.\endnote{Sister Lucia dos Santos, Third Memoir (1941); cf.\ \textit{Documenta\c{c}\~{a}o Cr\'{i}tica de F\'{a}tima: Sele\c{c}\~{a}o de documentos (1917--1930)} (F\'{a}tima: Santu\'{a}rio de F\'{a}tima, 2013).} From that afternoon passed into the universal Rosary of the Catholic faithful the brief decade prayer taught by Our Lady: ``O my Jesus, forgive us our sins, save us from the fires of hell, lead all souls to heaven, especially those most in need of Thy mercy.'' In Smyrna, Carthage, Antioch, Manresa, Avila, and Fatima alike, looking steadily at hell produces the same Christian posture: fearlessness before the world, vigilance over one's own soul, gratitude to the Redeemer, and sacrificial love for every neighbor still on the way.
